\documentclass{ximera}

\title{Limits at Infinity Activity Facilitator Guide}
\author{MATH 425: Calculus I}

\begin{document}
\begin{abstract}
    Working with the peers in your group, solve the following problems. Make sure to show and justify all your work. Make sure everyone in the group understands the solution and participates. Be prepared to report your answers to the whole class. 
\end{abstract}
\maketitle

\textbf{Student Learning Objectives:}
Students will be able to:
\begin{itemize}
    \item Evaluate limits at infinity using methods discussed in class.
    \item Create examples of functions with specific limit and discontinuity properties.
\end{itemize}

\textbf{Prerequisite Knowledge:} Students should be familiar with limits, end behavior and limits at infinity.  Students will also be required to understand types of discontinuities to explore how they may interact with limits.\\

This activity is more practice-based than other activities, as these limits are something students generally need additional support in order to fully understand.  Check with the instructor if they did not specify how much work students will be required to show for these limits.  The algebraic renormalization procedure may not be required for all cases, but it is necessary for the limits involving a square root in the division.  The final problem is more creative to get students to explore the interaction between limits and discontinuities.  If multiple groups get this far, have them trade examples and check each others' work.  


\begin{exercise}
  We frequently need to assess behaviour of functions involving $\lim_{x\rightarrow \infty}x^n$ and $\lim_{x\rightarrow -\infty}x^n$ (where $n$ is a real constant).  These particular limits often show up multiple times when anaylizing more complex algebraic functions.  Investiage these limits in each of the following circumstances. (Your group may find using graphs to visualize the behavior is helpful.)
  \begin{enumerate}
    \item $n$ is an even integer.\\
    \textcolor{blue}{If $n>0$,$\displaystyle\lim_{x\rightarrow \infty} x^n=\infty$ and $\displaystyle\lim_{x\rightarrow -\infty} x^n=\infty$\\
    If $n<0$, $\displaystyle \lim_{x\rightarrow \infty}x^n=0$ and $\displaystyle \lim_{x\rightarrow -\infty}x^n=0$}
    \item $n$ is an odd integer.\\
    \textcolor{blue}{If $n>0$, $\displaystyle\lim_{x\rightarrow \infty} x^n=\infty$ and $\displaystyle\lim_{x\rightarrow -\infty} x^n=-\infty$\\
    If $n<0$, $\displaystyle \lim_{x\rightarrow \infty}x^n=0$ and $\displaystyle \lim_{x\rightarrow -\infty}x^n=0$}
    \item $n$ is an odd integer.
    \item $n$ is rational with an even denominator. (You may need to test a few cases to determine the behavior here.)\\
    \textcolor{blue}{If $n>0$, then $\displaystyle \lim_{x\rightarrow \infty}x^n=\infty$, and if $n<0$, $\lim_{x\rightarrow \infty}x^n=0$.  For any case $\lim_{x\rightarrow -\infty}x^n$ is not defined.  The even root of a negative number is not a real value.}
    \item $n$ is rational with an odd denominator.\\
    \textcolor{blue}{If $n<0$, $\displaystyle \lim_{x\rightarrow \infty}x^n=0$ and $\displaystyle \lim_{x\rightarrow -\infty}x^n=0$\\
    If $n>0$ and the numerator of $n$ is even, $\displaystyle \lim_{x\rightarrow \infty}x^n=\infty$ and $\displaystyle \lim_{x\rightarrow -\infty}x^n=\infty$\\
    If $n>0$ and the numerator of $n$ is odd, $\displaystyle \lim_{x\rightarrow \infty}x^n=\infty$ and $\displaystyle \lim_{x\rightarrow -\infty}x^n=-\infty$}
  \end{enumerate}
  \textcolor{orange}{There are some tedious cases here, but the most important thing to remind students is that negative and positive integer exponents work differently.  If groups are missing cases remind them that the most important cases to check are going to be positive/negative and odd/even (for (d)).}
\end{exercise}

\begin{exercise}
    Evaluate the limits.\\
  \begin{enumerate}
    \item $\displaystyle\lim_{x\rightarrow -\infty}\frac{x^2-3x+2}{x-3}$\\
    \textcolor{blue}{$=\displaystyle \lim_{x\rightarrow -\infty} \frac{\frac{x^2}{x}-\frac{3x}{x}+\frac{2}{x}}{\frac{x}{x}-\frac{3}{x}}=\lim_{x\rightarrow -\infty} \frac{x-3+\frac{2}{x}}{1-\frac{3}{x}}=-\infty$}\\
    \textcolor{orange}{Check with the instructor to see if they will accept reasoning about degree of numerator and denominator for these problems.}
    
    \item $\displaystyle\lim_{x\rightarrow -\infty} (x^2-7x^5)$\\
    \textcolor{blue}{$=\displaystyle \lim_{x\rightarrow -\infty} x^2(1-7x^3)$\\
    $x^2\rightarrow \infty$ as $x\rightarrow -\infty$ and $1-7x^3\rightarrow \infty$ as $x\rightarrow \infty$, so $\displaystyle\lim_{x\rightarrow -\infty} (x^2-7x^5)=\infty$}
    
    \item $\displaystyle\lim_{s\rightarrow \infty}\frac{\sqrt{1+4s^6}}{2-s^3}$\\
    \textcolor{blue}{$\displaystyle\lim_{s\rightarrow \infty}\frac{\sqrt{s^6(\frac{1}{s^6}+4)}}{2-s^3}$\\
    $=\displaystyle \lim_{x\rightarrow \infty}\frac{\frac{s^3}{s^3}\sqrt{\frac{1}{s^6}+4}}{\frac{2}{s^3}-\frac{s^3}{s^3}}=\lim_{x\rightarrow \infty}\frac{\sqrt{\frac{1}{s^6}+4}}{\frac{2}{s^3}-1}=\frac{\sqrt{4}}{-1}=-2$}\\
    \textcolor{orange}{Students may use powers here to argue about the limit, but they should also be showing the algebraic renomalization procedure for problems with roots.}
    
    \item $\displaystyle\lim_{x\rightarrow \infty}(\sqrt{9x^2+x}-3x)$ (Hint: this may benefit from multiplication by the conjugate over itself.)\\
    \textcolor{blue}{$\displaystyle \lim_{x\rightarrow \infty}\sqrt{9x^2+x}-3x\bigg(\frac{\sqrt{9x^2+x}+3x}{\sqrt{9x^2+x}+3x}\bigg)$\\
    $=\displaystyle \lim_{x\rightarrow \infty} \frac{9x^2+x-9x^2}{\sqrt{9x^2+x}+3x}=\lim_{x\rightarrow \infty}\frac{x}{|3x| \sqrt{1+\frac{1}{9x}}+3x}=\lim_{x\rightarrow \infty}\frac{x}{x(3\sqrt{1+\frac{1}{9x}}+3)}=\lim_{x\rightarrow \infty}\frac{1}{3\sqrt{1+\frac{1}{9x}}+3}=\frac{1}{6}$}

    \item $\displaystyle \lim_{x\rightarrow \infty} \sqrt{3+x^2}-x$ (Hint: this may benefit from multiplication by the conjugate over itself.)\\
    \textcolor{blue}{Multiply by the conjugate over itself: $\displaystyle \lim_{x\rightarrow \infty}\sqrt{3+x^2}-x\bigg(\frac{\sqrt{3+x^2}+x}{\sqrt{3+x^2}+x}\bigg)$\\
    $=\displaystyle \lim_{x\rightarrow \infty} \frac{3+x^2-x^2}{\sqrt{3+x^2}+x}=\lim_{x\rightarrow \infty}\frac{3}{x\sqrt{\frac{3}{x^2}+1}+x}=\lim_{x\rightarrow \infty}\frac{\frac3x}{\frac{x}{x}\sqrt{\frac{3}{x^2}+1}+\frac{x}{x}}$\\
    $\displaystyle=\lim_{x\rightarrow \infty}\frac{\frac3x}
    {\sqrt{\frac{3}{x^2}+1}+1}=\frac{0}{2}=0$}

    \item $\displaystyle\lim_{x\rightarrow \infty}\frac{1+x^4}{x^6+1}$\\
    \textcolor{blue}{$\displaystyle =\lim_{x\rightarrow \infty}\frac{\frac{1}{x^6}+\frac{x^4}{x^6}}{\frac{x^6}{x^6}+\frac{1}{x^6}}=\frac{0}{1}=0$}
    
    \item $\displaystyle \lim_{x\rightarrow \infty} 1.5^{-x}-7$\\
    \textcolor{blue}{$\displaystyle =\lim_{x\rightarrow \infty} \frac{1}{1.5^x}-7=0-7=-7$}

    \item $\displaystyle \lim_{x\rightarrow -\infty} 1.5^{-x}-7$\\
    \textcolor{blue}{$\displaystyle =\lim_{x\rightarrow -\infty} \frac{1}{1.5^x}-7=\infty$}

  \end{enumerate}
\end{exercise} 
    

\begin{exercise}
    

 Create an example of a function which satisfies the following.  Trade your examples with a classmate or another group and challenge them to figure out which function fits each problem!\\
 \textcolor{orange}{There are infinitely many possibilities for all of these problems. When interacting with students on these create an example problems, ask students to defend their choices.  This gives you the opportunity to check their answer, and also makes sure they understand the example they chose.}
 \begin{enumerate}
    \item A function $f$ with no infinite discontinuities where $\displaystyle\lim_{x\rightarrow \infty}f(x)=3$. \\
    \textcolor{blue}{ An example would be $f(x)=\frac{3x^2-3}{x^2+1}$}

    \item A function $f$ with an infinite discontinuity at $x=3$ and a removable discontinuity at $x=1$.\\
    \textcolor{blue}{An example would be $f(x)=\frac{x-1}{(x-1)(x-3)}$}
    
    \item A function with $\displaystyle\lim_{x\rightarrow \infty}f(x)=-\infty$ and $\displaystyle\lim_{x\rightarrow -\infty}f(x)=\infty$.\\
    \textcolor{blue}{ An example would be $f(x)=-x^3$}
 \end{enumerate}

\end{exercise}


\end{document}